% Source: https://artofproblemsolving.com/community/c3h2793114_2004__midmichigan_mathematical_olympiad___grades_56
% Stable identity: aops-topic-2793114-q03
Whenever possible, group two 1-cent coins into one 2-cent unit, five 2-cent units into one 10-cent unit, and two 10-cent units into one 20-cent unit. A subset of grouped units corresponds to a subset of original coins, so the grouped collection still cannot make 100. It has at most one 1, four 2 s, one 10, four 20 s, and one 50: otherwise five 20 s or two 50 s make 100. The small coins total at most 9. If both 50 and 10 occur, at most one 20 may occur, since $50+10+20+20=100$; the total is then at most 89. If 50 is absent, the total is at most 99. If 50 occurs but 10 does not, the total is at most $50+80+9=139$. This is attained by one 50, four 20 s, and nine 1 s. Without 50 a subset totals at most 89; with 50 the other coins cannot total 50, since they contribute a multiple of 20 plus at most 9. Thus $\boxed{139\text{ cents}}$.
